\section{下降}\srcpage{15}

我们先回忆两个一般结果。对任意概形（或代数叠）$S$ 与任意 $K \in \Dbc(S)$，记
\[
  \chi_K \colon S \to \Z
\]
为 $K$ 的 Euler--Poincar\'e 示性函数，它由
\[
  \chi_K(s) = \sum_i (-1)^i \dim_{\C} H^i(K_s), \qquad \forall s \in S,
\]
定义，其中 $K_s$ 是 $K$ 在 $s$ 处的纤维。

\begin{lemma}[Deligne, {[Il]} 推论 2.10]
设 $f \colon Y \to Z$ 是概形的紧合态射，甚至是代数叠的可表紧合态射。那么，若
$K, K' \in \Dbc(Y)$ 在 $Y$ 上局部同构，则复形 $f_*K$ 与 $f_*K'$ 在 $Z$ 上有相同的
Euler--Poincar\'e 示性函数 $\chi_{f_*K} = \chi_{f_*K'}$。
\end{lemma}

\begin{lemma}
设 $S$ 是概形，甚至是代数叠；设 $V \to S$ 是常秩 $r$ 的向量丛（$r$ 有限），它定义射影丛
$\pi \colon P = \Pv(V) \to S$；设 $K$ 是 $P$ 上的不可约反常层。那么，存在 $S$ 上的反常层
（必然不可约）$L$，使得 $K$ 同构于 $\pi^*(L)[r-1]$，当且仅当
Euler--Poincar\'e 示性函数 $\chi_K$ 沿 $\pi$ 的纤维为常数。
\end{lemma}

\begin{proof}
「仅当」部分平凡。因此设 $\chi_K$ 沿 $\pi$ 的纤维为常数。由不可约反常层的结构定理，有
$K \cong i_* j_{!*} E[\dim Q]$，其中 $i \colon Q \hookrightarrow P$ 是 $K$ 的不可约闭支撑，
$j \colon Q^\circ \hookrightarrow Q$ 是稠密开集，$E$ 是 $Q^\circ$ 上的不可约局部系统。

首先验证 $Q = \pi^{-1}(T)$，其中 $T$ 是 $S$ 的不可约闭子集。其次验证可取 $Q^\circ$ 形如
$\pi^{-1}(T^\circ)$，其中 $T^\circ$ 是 $T$ 的稠密开集。最后验证 $E$ 来自 $T^\circ$ 上的
不可约局部系统 $F$，因为 $\pi$ 的纤维单连通。

所求的 $S$ 上反常层 $L$，就是 $F[\dim T]$ 到 $T$ 的中间延拓在整个 $S$ 上的零延拓。
\end{proof}

现在研究 $\Aut'_E$ 到 $\Fib_n(X)$ 上反常层 $\Aut_E$ 的下降。为此，考虑 Drinfeld 引入的叠
$\cY^d_n$，它分类满足下述条件的二元组 $(\cL, s^\bullet)$：
\srcpage{16}$\cL$ 是次数为 $n(n-1)(g-1)+d$ 的向量丛，而
$s^\bullet = (s^1, \dots, s^n)$ 是一列 $\cO_X$-模单射
\[
\begin{array}{rcl}
(\OmX)^{\otimes n-1} & \xhookrightarrow{s^1} & \cL \\
 & \cdots & \\
(\OmX)^{\otimes (n-1)+\cdots+(n-i)} & \xhookrightarrow{s^i} & \bigwedge^i \cL \\
 & \cdots & \\
(\OmX)^{\otimes \frac{n(n-1)}{2}} & \xhookrightarrow{s^n} & \bigwedge^n \cL
\end{array}
\]
它们满足 Pl\"ucker 关系，使得 $s^\bullet$ 在 $X$ 的一般点处给出 $\cL$ 纤维的一个
子空间完全旗。

叠 $\cY^d_n$ 是三元组 $(\cL, \cL_\bullet, s_\bullet)$ 的叠的一个
「部分紧化」，其中
\[
  \cL_\bullet = ((0) = \cL_0 \subset \cL_1 \subset \cdots \subset \cL_n)
\]
是 $\cL$ 的 $\cO_X$-子模旗，而 $s_\bullet$ 是一列 $\cO_X$-模同构
$s_j \colon (\OmX)^{\otimes j-1} \xrightarrow{\sim} \cL_{n-j+1}/\cL_{n-j}$，
$j = 1, \dots, n$。事实上，由
\[
\begin{aligned}
  s^i = s_n \otimes s_{n-1} \otimes \cdots \otimes s_{n-i+1} \colon {}
  & (\OmX)^{\otimes (n-1)+(n-2)+\cdots+(n-i)} \\
  & \xrightarrow{\sim}
  \cL_1 \otimes (\cL_2/\cL_1) \otimes \cdots \otimes
  (\cL_i/\cL_{i-1}) \cong \bigwedge^i \cL_i
\end{aligned}
\]
定义的态射
\[
  (\cL, \cL_\bullet, s_\bullet) \mapsto (\cL, s^\bullet)
\]
把这个三元组叠与开集 $\cY^{d\circ}_n \subset \cY^d_n$ 等同，该开集由满足
下述条件的 $(\cL, s^\bullet)$ 组成：当 $i = 1, \dots, n-1$ 时，单射
\[
  s^i \colon (\OmX)^{\otimes (n-1)+\cdots+(n-i)} \hookrightarrow \bigwedge^i \cL
\]
的余核无挠。

环面 $\GG_m^n$ 按下式作用在 $\cY^d_n$ 上：
\[
  (t_1, \dots, t_n) \cdot (\cL, (s^1, \dots, s^n)) =
  (\cL, (t_1 s^1, \dots, t_n s^n)).
\]

这个作用保持开集 $\cY^{d\circ}_n$，它在三元组叠上由
\[
  (t_1, \dots, t_n) \cdot (\cL, \cL_\bullet, (s_1, \dots, s_{n-1}, s_n)) =
  (\cL, \cL_\bullet,
  (t_n t_{n-1}^{-1} s_1, \dots, t_2 t_1^{-1} s_{n-1}, t_1 s_n))
\]
给出。

特别地，第 4 节引入的代数叠 $\cU_n$（相应地，
$\overline{\cU}_n = [\cU_n/\GG_m]$），可以等同于
$\cY^{0\circ}_n$ 对 $\GG_m$（相应地，$\GG_m^2$）作用的商。该作用经由嵌入
$\GG_m \hookrightarrow \GG_m^n$，$t \mapsto (t, t^2, \dots, t^n)$（相应地，
$\GG_m^2 \hookrightarrow \GG_m^n$，
$(t, t') \mapsto (t, t^2t', t^3t'^3, \dots, t^n t'^{\frac{n(n-1)}{2}})$）给出。

\srcpage{17}
\begin{proposition}[{[B-G]} 命题 1.2.2]
遗忘态射 $(\cL, s^\bullet) \to (\cL, s^1)$ 过商为代数叠的态射
\[
  \tilde{f} \colon \widetilde{\cY}^d_n := [\cY^d_n/\GG_m^n]
  \to {\Fib'}^{n(n-1)(g-1)+d}_n(X),
\]
它是可表且紧合的。
\end{proposition}

考虑第 4 节引入的 $\overline{\cU}_n = [\cY^{0\circ}_n/\GG_m^2]$ 上的复形
$(\overline{g}_n)^*\Psi$，并记 $\Phi \in \Dbc(\widetilde{\cY}^0_n)$ 为它沿复合
态射的紧支撑前推；该复合态射由商映射
$[\cY^{0\circ}_n/\GG_m^2] \to [\cY^{0\circ}_n/\GG_m^n]$ 与包含
$[\cY^{0\circ}_n/\GG_m^n] \subset [\cY^0_n/\GG_m^n] =
\widetilde{\cY}^d_n$ 组成。\yizhu{此处原文印作
$\widetilde{\cY}^d_n$。按上式末端的等号，$[\cY^0_n/\GG_m^n]$ 应等于
$\widetilde{\cY}^0_n$（见 p17 处的定义）；且 $\Phi$ 本身是
$\widetilde{\cY}^0_n$ 上的复形。故 $d$ 应为 $0$。译文照原文保留。}{原文笔误，I（原文 p15 定义互证）}

对 $X$ 上任意秩 $n$ 局部系统 $E$，有函子 $M^d_{n,E}$ 的一个变体
\[
  N^d_{n,E} \colon \Dbc(\widetilde{\cY}^0_n) \to \Dbc(\widetilde{\cY}^d_n),
  \quad K \mapsto p_*p'^*(K \otimes q^*\cL^d_E)[nd],
\]
它借助图表
\[
\begin{tikzcd}
 & \widetilde{\cZ}^d_n \arrow[dl, "p'"'] \arrow[dr, "p"] \arrow[r, "q"]
 & \Coh^d_0(X) \\
 \widetilde{\cY}^0_n & & \widetilde{\cY}^d_n
\end{tikzcd}
\]
定义。其中
\[
  \widetilde{\cZ}^d_n = \widetilde{\cY}^0_n
  \times_{\Fib_n(X),p'} \Mod^d_n(X)
\]
是三元组 $(\cL', s'^\bullet, \cL' \subset \cL)$ 的代数叠对
$\GG_m^n$ 的商，这里 $(\cL', s'^\bullet) \in \cY^0_n$ 且
$(\cL' \subset \cL) \in \Mod^d_i(X)$\yizhu{原文此处印作 $\Mod^d_i(X)$；按本页纤维积的定义与 $\cL'$ 的秩 $n$，下标应为 $n$。译文照原文保留。}{原文笔误，I（p17 上下文推断，无外部依据）}；而 $p$、$p'$、$q$ 由下述箭头诱导：
这些箭头把这样的三元组分别送到 $(\cL, s^\bullet)$、$(\cL', s'^\bullet)$ 与
$\cL/\cL'$，其中 $s^i$ 是 $s'^i$ 与包含
$\bigwedge^i \cL' \subset \bigwedge^i \cL$ 的复合。

\begin{lemma}
设 $E$ 是 $X$ 上秩 $n$ 的局部系统。对任意整数 $d \ge 0$，有
\[
  \Aut'_E\,|\,{\Fib'}^{n(n-1)(g-1)+d}_n = \tilde{f}_*N^d_{n,E}(\Phi).
\]
\end{lemma}

为证明 $\Aut'_E$ 沿遗忘截面的态射 $\Fib'_n(X) \to \Fib_n(X)$ 下降，Frenkel、Gaitsgory
与 Vilonen 用引理 7.4、命题 7.3 与引理 7.1 证明 $\Aut'_E$ 的
Euler--Poincar\'e 示性函数不依赖 $E$。然后他们在 $E$ 为秩 $n$ 平凡局部系统时计算这个
Euler--Poincar\'e 示性函数，最后用引理 7.2 得出结论。
